% ==========================================================================
% INFO 4617 Web Data Science -- Week 7 Friday activity: Issue to pull request
% Build: `make` in handouts/, or `latexmk -pdf issue-to-pull-request.tex`
% here. The figures are annotated from real captures (img/IMAGES.md). A
% screen that needs a signed-in browser is a hand capture: \handcapture
% prints it once img/NAME.png exists, and leaves it out until then.
% ==========================================================================
\makeatletter\def\input@path{{../common/}{handouts/common/}}\makeatother
\documentclass{handout}
\usepackage{amssymb}   % \checkmark

\title[INFO 4617 · Week 7 · Friday activities · Issue to pull request]{Issue to pull request}
\subtitle{Week 7 Friday activity · hard · about 30 minutes}
\author{Brian C. Keegan, Ph.D.}
\institute{Department of Information Science, University of Colorado Boulder}
\date{October 2, 2026}

% A numbered figure: [width]{path from the repo root}{alt text}{caption}
\newcounter{handoutfig}
\newcommand{\activityfigure}[4][\linewidth]{%
  \par\smallskip\refstepcounter{handoutfig}%
  {\centering
   \BeginAccSupp{method=pdfstringdef,Alt={#3}}%
   \includegraphics[width=#1]{#2}%
   \EndAccSupp{}\par}%
  \smallskip{\small\raggedright\textbf{Figure \thehandoutfig.} #4\par}}
% A hand capture: {file in handouts/week-07/img/}{alt text}{caption}
\newcommand{\handcapture}[3]{%
  \IfFileExists{handouts/week-07/img/#1}%
    {\activityfigure{handouts/week-07/img/#1}{#2}{#3}}%
    {\IfFileExists{../../handouts/week-07/img/#1}%
      {\activityfigure{handouts/week-07/img/#1}{#2}{#3}}{}}}

\setlist[enumerate]{leftmargin=*,itemsep=1.5pt,topsep=2pt}
\setlist[itemize]{leftmargin=*,itemsep=1pt,topsep=2pt}

\begin{document}
\maketitle

\begin{frame}{What you do}
  \begin{columns}[T]
    \begin{column}{0.57\textwidth}
      You fix the problem that an issue describes, in GitHub's web editor.
      Then you open a pull request that closes the issue when it merges. You
      do all of it in the browser: you need no terminal and no copy of the
      book on your laptop.

      Your chapter's board in week 7's slides lists issues that have no pull
      request. Start with one that a triage comment calls \textbf{small}.

      Work in pairs. One partner works in the browser, and the other reads
      this handout aloud.
    \end{column}
    \begin{column}{0.39\textwidth}
      \begin{block}{The three Friday activities}
        \small
        \begin{tabular}{@{}ll@{}}
          Triage an issue & easy\\
          Review a pull request & medium\\
          \textbf{Issue to pull request} & \textbf{hard}
        \end{tabular}
      \end{block}
      \begin{alertblock}{Do not}
        \small
        \begin{itemize}
          \item click \textbf{Merge pull request};
          \item change a file in \texttt{notebooks/} or \texttt{images/}: a script makes the notebooks from the chapters;
          \item put two fixes in one pull request.
        \end{itemize}
        The instructor merges, closes, and resolves conflicts.
      \end{alertblock}
    \end{column}
  \end{columns}
\end{frame}

\begin{frame}{1 · Claim the issue}
  \begin{enumerate}
    \item Sign in to GitHub.
    \item Open the issue
      (\url{https://github.com/cuinfoscience/Web-Data-Science-Book/issues/}
      and its number). Read it, and read its comments.
    \item Write the comment \texttt{I'm taking this.} and click
      \textbf{Comment}. Then no other pair starts the same fix.
  \end{enumerate}
\end{frame}

\begin{frame}{2 · Open the chapter in the editor}
  \begin{enumerate}
    \item Open the book at the chapter:
      \url{https://cuinfoscience.github.io/Web-Data-Science-Book/}
    \item Click \textbf{Edit this page} ①, at the bottom of the right-hand
      sidebar (Figure~\ref{fig:edit}). GitHub opens the chapter's source
      file, \texttt{ch-NN-name.qmd}.
    \item Click the pencil icon, \textbf{Edit this file}. If GitHub asks you
      to \textbf{fork} the repository, accept. A fork is your own copy of the
      book, and GitHub makes it for you. The steps after this are the same.
  \end{enumerate}
  \activityfigure[0.64\linewidth]{handouts/week-07/img/edit-this-page_annotated.pdf}%
    {Screenshot of chapter 4 of the book's website. Marker 1 points to Edit
     this page, at the bottom of the right-hand sidebar.}%
    {A chapter of the book (here chapter 4). \textbf{Edit this page} ① is at
     the bottom of the right-hand sidebar, above \textbf{Report an issue}.}%
  \label{fig:edit}
\end{frame}

\begin{frame}{3 · Make one change}
  \begin{enumerate}
    \item Click in the editor and press \texttt{Ctrl+F} (\texttt{Cmd+F} on a
      Mac). This opens the editor's own search bar ① (Figure~\ref{fig:search}).
      Your browser's search cannot see the whole file. Search for words from
      the section, such as its heading ②.
  \end{enumerate}
  \activityfigure[0.78\linewidth]{handouts/week-07/img/editor-search_annotated.pdf}%
    {Screenshot of GitHub's web editor with its search bar open. Marker 1 is
     on the search box, which holds "missing manual". Marker 2 is beside the
     heading the search found.}%
    {GitHub's web editor, with its own search bar open ①. The search for
     \textit{missing manual} found the callout's heading ②.}%
  \label{fig:search}
  \begin{enumerate}[resume]
    \item Make the change that the issue asks for. Change nothing else.
    \item Keep the Markdown correct. Put an empty line above a heading
      (\texttt{\#\#}) or a callout (\texttt{:::}). Write a link as
      \breakcode{[text](https://...)}, with no space between \texttt{]} and
      \texttt{(}.
    \item Click \textbf{Preview} ① (Figure~\ref{fig:preview}). Compare the
      old line ② with the new line ③. Then click \textbf{Commit changes…} ④.
  \end{enumerate}
  \activityfigure[0.78\linewidth]{handouts/week-07/img/preview-diff_annotated.pdf}%
    {Screenshot of the editor's Preview tab after an edit. Marker 1 is on the
     Preview tab. Marker 2 is on the old line, in red. Marker 3 is on the new
     line, in green. Marker 4 is under the Commit changes button.}%
    {The \textbf{Preview} tab ① after an edit. The old line ② is red, and
     the new line ③ is green. \textbf{Commit changes…} ④ is at the top
     right.}%
  \label{fig:preview}
\end{frame}

\begin{frame}{4 · Commit to a new branch}
  \begin{enumerate}
    \item In \textbf{Commit message} ①, say what you changed and where, for
      example \texttt{Ch. 5: say where requests stores redirect hops}.
    \item Select \textbf{Create a new branch for this commit and start a pull
      request} ② (Figure~\ref{fig:commit}). The dialog selects
      \textit{Commit directly to the main branch} first: change it.
      \texttt{main} refuses direct commits.
    \item Click the green button ③. When a new branch is selected, the
      button says \textbf{Propose changes}. (In a fork, GitHub makes the
      branch for you, and the button always says \textbf{Propose changes}.)
  \end{enumerate}
  \activityfigure[0.66\linewidth]{handouts/week-07/img/commit-dialog_annotated.pdf}%
    {Screenshot of the Commit changes dialog. Marker 1 is beside the commit
     message. Marker 2 is beside the option Create a new branch for this
     commit and start a pull request. Marker 3 is beside the green button.}%
    {The \textbf{Commit changes} dialog: the commit message ①, the option
     that makes a new branch and starts a pull request ②, and the green
     button ③, which says \textbf{Propose changes} once that option is
     selected.}%
  \label{fig:commit}
\end{frame}

\begin{frame}{5 · Open the pull request}
  \begin{enumerate}
    \item GitHub shows \textbf{Comparing changes}: your branch and
      \texttt{main}, with the difference between them. Make sure that only
      your change shows. Click \textbf{Create pull request}.
    \item Write a title that says what and where. Do not write
      \texttt{Update ch-05-protocols.qmd}. Write \texttt{Ch. 5: say where
      requests stores redirect hops}.
    \item The description starts with the book's template. Replace each
      \texttt{[bracketed hint]}:
      \begin{itemize}
        \item \texttt{Closes \#N}: the issue's number, for example
          \texttt{Closes \#175}. GitHub closes the issue when the pull
          request merges.
        \item \textbf{Location}, \textbf{Problem}, \textbf{Why}, and
          \textbf{Change}: the file and the heading, what was wrong, who it
          affects, and what you changed.
        \item \textbf{AI assistance}: \texttt{none}, or the tool and what it
          did.
        \item \textbf{Checks}: tick the first box. After the second box,
          write \texttt{Edited in the browser, so notebooks not
          regenerated.} Tick the third box when the Render check is green.
      \end{itemize}
    \item Click \textbf{Create pull request}.
  \end{enumerate}
  \handcapture{pr-template.png}{Screenshot of the Open a pull request form
    with the book's template in the description box.}{The \textbf{Open a pull
    request} form, with the book's template in the description, before
    \textbf{Create pull request}.}
\end{frame}

\begin{frame}{6 · After you open it}
  Then watch your pull request. If you see one of these, do what the
  right-hand column says.

  \smallskip
  {\small
  \begin{tabular}{@{}>{\raggedright\arraybackslash}p{0.3\linewidth}%
                  >{\raggedright\arraybackslash}p{0.66\linewidth}@{}}
    \toprule
    \textbf{You see} & \textbf{Do this}\\
    \midrule
    \textbf{Render} with a red \texttimes &
      Click \textbf{Details}. Find the first \texttt{ERROR} line. A broken
      cross-reference (\texttt{@sec-...}) or a code block that does not
      close is the usual cause. Fix it on the same branch, as in the next
      row.\\[3pt]
    A review that asks for changes &
      Open your pull request, and click \textbf{Files changed}. Click the
      file's \textbf{$\cdots$} menu, then \textbf{Edit file}. Make the change, and
      commit it to the same branch. Do not open a new pull request.\\[3pt]
    \textbf{Notebook sync} with a red \texttimes &
      Normal for a change made in the browser. The instructor regenerates the
      notebooks after merging.\\[3pt]
    \textit{This branch has conflicts that must be resolved} &
      Leave it. The instructor resolves conflicts.\\
    \bottomrule
  \end{tabular}\par}

  \handcapture{pr-edit-file.png}{Screenshot of the Files changed tab of a
    pull request with the file's menu open on Edit file.}{\textbf{Files
    changed} on your own pull request: the file's \textbf{$\cdots$} menu, open on
    \textbf{Edit file}.}

  \medskip
  \textbf{Next.} Ask a pair in your chapter to review your pull request
  (\breakcode{review-a-pull-request.pdf}), and review theirs.
\end{frame}
\end{document}
